% Generated draft: requires mathematical and source review.
\documentclass[11pt]{article}
\usepackage[margin=25mm]{geometry}
\usepackage{fontspec}
\setmainfont{Latin Modern Roman}
\usepackage{amsmath,amssymb,mathtools,unicode-math}
\setmathfont{Latin Modern Math}
\usepackage{enumitem,needspace,xurl,hyperref}
\setlength{\parindent}{0pt}
\setlength{\parskip}{5pt}
\setlength{\emergencystretch}{3em}
\newcommand{\N}{\mathbb{N}}
\newcommand{\Z}{\mathbb{Z}}
\newcommand{\Q}{\mathbb{Q}}
\newcommand{\R}{\mathbb{R}}
\newcommand{\C}{\mathbb{C}}
\newcommand{\F}{\mathbb{F}}
\newcommand{\A}{\mathbb{A}}
\newcommand{\PP}{\mathbb{P}}
\newcommand{\E}{\mathbb{E}}
\newcommand{\eps}{\varepsilon}
\newcommand{\dd}{\,\mathrm{d}}
\begin{document}
\section*{UnsolvedMath ID: 6000018}
\textbf{Source problem code:} catalogue entry
\subsection{Definitions and mathematical statement}
Let $(X,\mathcal A,\lambda)$ be a sample space with a reference measure and let $\Theta$ be a parameter space. A statistical family of densities is a collection $f(x;\theta)\geq0$ satisfying
\[
\int_X f(x;\theta)\,d\lambda(x)=1\qquad(\theta\in\Theta).
\]
The pairs $(x,\theta)$ form the sample--parameter product $X\times\Theta$. The family gives a conditional distribution on $X$ for each parameter, not a joint probability measure on the product unless additional parameter weighting is supplied.

Construct and study a useful geometry on $X\times\Theta$ derived from the statistical family, while accounting for the selection of a statistically meaningful coordinate system on $\Theta$. If $\theta=\varphi(\eta)$ is a change of parameter coordinates, the same family is represented as $\widetilde f(x;\eta)=f(x;\varphi(\eta))$; a proposed construction should explain its relation to this reparametrization and the statistical criterion for any preferred coordinates.

The source does not prescribe a metric, connection, divergence, or other geometric structure; nor does it define usefulness or statistical meaningfulness. It also supplies no smooth structure on the sample space, regularity assumptions on the family, or joint measure on the product. Thus the target is a geometry-construction research program with criteria still to be specified, not a unique existence theorem for a named geometric object.
\subsection{Short English statement}
Develop geometry on the space of sample--parameter pairs for a statistical model, and explain how statistically meaningful parameter coordinates enter the construction. The desired structure and its criteria of usefulness are unspecified.
\subsection{Sources}
\begin{itemize}
\item Furuhata, Matsuzoe, and Urakawa, \emph{Open Problems in Affine Differential Geometry and Related Topics} (1998), item 6(a), scanned PDF pages 1--2, \url{https://www.jstage.jst.go.jp/article/iis/4/2/4_2_125/_pdf/-char/en}; accessed 2026-07-29, manual OCR from page images recorded in the provenance. The record attributes the question to Kwoichi Tandai.
\item The supplied research summary mentions density-space geometry, optimal-transport-based information geometry, and sample-space geometry as partial tools, without individual publication citations.
\item Local registry \texttt{data/unsolvedmath/problems.json}, permanent UnsolvedMath ID $6000018$, source code catalogue entry, \emph{Geometry on Sample-Parameter Product Spaces}. Dataset attribution: Ulam AI and the UnsolvedMath Contributors.
\item The supplied rights note says the J-STAGE article page marks this article CC BY 4.0.
\end{itemize}
\end{document}
